\clearpage
\fancyhead[L]{\small CONNAISSANCE DES TEMS.}
\hypertarget{native-559-lower}{}\label{native-559-lower}
\begin{center}\small Göttingische gelehrte Anzeigen. No. 8. Pages 873..876. 3 June 1813.\\[3pt]\rule{.72\linewidth}{.3pt}\end{center}

\emph{Connaissance des tems ou des mouvemens célestes, à l'usage des astronomes et des navigateurs, pour l'an} 1813, publiée par le bureau des longitudes. Juillet 1811. Paris. 228 pages in octavo.

\emph{Connaissance des tems etc.} pour l'an 1814. Avril 1812. 240 pages.

The table of geographical positions of the principal places on Earth has again grown considerably, especially through drawing on positions determined during the voyages of \textsc{Dentrecasteaux}, \textsc{Humboldt}, and \textsc{Krusenstern}. In the volume for 1814 the table thus contains over 1600 entries, and consequently over 100 more than the volume for 1812. The \emph{supplements} to these two volumes are few this time. The 1813 volume contains only two short papers by \textsc{Laplace}. In the first, he makes an important addition to the theory of the so-called method of least squares. By an ingenious analysis he shows that, with the combination of the basic equations prescribed by this method, the mean error to be expected in the results is minimized as soon as the number of basic equations is very large, whatever the probability law of the errors may be. This investigation is carried out in greater detail in the author's larger work on the calculus of probabilities, which appeared recently. The other supplement concerns the long-period equation of the Moon. Its necessity has been inferred from the discrepancy among the epochs of the Moon's mean longitude obtained for 1692, 1756, 1779, and 1801, \SourceLeaf{560}although it has yet to be established entirely satisfactorily in the theory. In \textsc{Bürg}'s lunar tables it was taken to be proportional to the product of the sine of the quantity obtained by taking twice the longitude of the lunar node plus the longitude of the lunar perigee less three times the longitude of the solar perigee, and its coefficient was taken as $= -14''$. \textsc{Laplace} now announces that he has stronger reasons (which he does not explain further) for taking it to be proportional to the cosine of the sum of twice the longitude of the lunar node and the longitude of the lunar perigee, and sets its coefficient $= -13''96$. (In his new lunar tables, \textsc{Burckhardt}, as we already remarked in our review of them. took the same form as his basis and assumed the coefficient $= -12''5$.) We eagerly await the author's fuller account of his theoretical investigations of this delicate question.

The supplements to the 1814 volume are as follows. 1) On the origin of comets, by \textsc{Lagrange}. The great geometer (since taken from the sciences by death, at an advanced age yet still too soon) extends \textsc{Olbers}'s well-known hypothesis on the origin of the new planets to comets. With the elegance we had always come to expect from him, he works out the conditions under which a fragment thrown from a planet moving in a circular orbit describes a parabolic orbit. As one can easily see (\textsc{Lagrange}, however, did not state this explicitly), the limits of the relative speed at which the nascent comet is propelled are $\BareRoot{2}-1$ and $\BareRoot{2}+1$, taking the speed of the planet as the unit. What is remarkable, however, is the proposition put forward here by \textsc{Lagrange} (likewise easy to prove) that the speed $\BareRoot{3}$ marks the dividing line between prograde and retrograde comets, so that a prograde comet (relative to the planet's orbit) requires a lower speed and a retrograde one a higher speed. Suppose a planet at a distance a hundred times as great as that of our Earth were shattered into several pieces by the action of an elastic fluid suddenly released inside it, from whatever cause. The explosion would only need to be strong enough to impart a speed twelve or fifteen times as great as that of a 24-pound cannonball for elliptical or parabolic comets of all possible dimensions to be \emph{produced in all possible directions}. We must, however, disagree with this latter claim. Only comets whose line of nodes in the former planetary orbit very nearly coincides with the line of apsides could arise in this way if they are to come down as far as the region of Earth's orbit at their closest approach to the Sun. Yet only a small proportion of all the cometary orbits computed so far are of a kind for which one could grant the possibility of such an origin. Practical astronomers will in any case hardly be inclined to see fragments of solid planetary bodies in comets, whose whole appearance points to an entirely different physical constitution. 2) Some notes from \textsc{Krusenstern}'s voyage around the world. 3) A table arranged to make it very easy to convert sidereal time to mean solar time, by \textsc{Burckhardt}. 4) An extract from the meteorological observations made in Paris in 1809. As usual, both volumes end with the list of members of the Bureau of Longitudes.
